% Dient zur Konfiguration von Layout und LaTeX
% Befehle zur Ausgabe von Tactical Training Team
\newcommand{\TTT}{Tactical Training Team}

% Zeilenabstand
\linespread{1.15}
\setlength{\parindent}{0pt} 

%Seitenränder
\geometry{
	left=30mm,
	right=20mm,
	top=30mm,
	bottom=30mm,
}

%Verzeichnisebenen Bilder und Tabellen
\counterwithin{figure}{section}
\counterwithin{table}{section} 

% Absand Aufzählung
\setlist[1]{itemsep=0pt}

% Farben
\usepackage{xcolor}
\definecolor{blue}{RGB}{0,0,120}
\definecolor{green}{RGB}{0,153,0}
\definecolor{backcolor}{rgb}{0.9,0.9,0.9}
\definecolor{yellow}{RGB}{255,210,0}
\definecolor{gold}{RGB}{218,165,32}
\definecolor{DarkOrchid}{RGB}{92,30,123}

% Url-Einstellungen
\hypersetup{%
	hidelinks,
	pdfauthor={\TTT}
}

% Einstellungen Kopf- und Fußzeile
\pagestyle{empty}
\clearscrheadfoot{}
\ofoot[]{\textbf{\vrule}\,\pagemark}
%\ihead[]{Tactical Training Team}
%\ohead[]{\headmark}
\chead[]{\includegraphics[width=1.0\textwidth]{../img/header}}
\addtolength{\headheight}{10mm}
\addtolength{\footskip}{-10mm}
%\setheadsepline[\textwidth]{1pt}[\color{800000}]

% Ebenen Nummerierungen
\setcounter{secnumdepth}{4} %Nummerierungsebenen
\setcounter{tocdepth}{1}   %Ebenen Inhaltsverzeichnis

%Neuer Splatentyp für Tabellen (dient dazu, dass die erste Splate zentiert werden kann)
\newcolumntype{C}[1]{>{\centering\arraybackslash}m{#1}}
%linksbündiger Spaltentyp
\newcolumntype{P}[1]{>{\raggedright\arraybackslash}p{#1}}

%Eigenen Befehl für Newline definieren
\newcommand{\nl}{\newline}

%Autoref in Schreibmaschinenschrift
%\let\oldautoref\autoref
%\renewcommand{\autoref}[1]{\texttt{\oldautoref{#1}}}

%Nameref in Schreibmaschinenschrift und mit >> <<
\makeatletter
\shorthandon{"}
\AtBeginDocument{%
	\LetLtxMacro\oldnameref\nameref
	\DeclareRobustCommand*{\nameref}{%
		\@ifstar{\my@nameref*}{\my@nameref{}}%
	}%
	\newcommand*{\my@nameref}[2]{%
		\texttt{>>\oldnameref#1{#2}<<}%
	}%
}
\shorthandoff{"}
\makeatother

\renewcommand*{\familydefault}{\sfdefault}	
%\renewcommand*\chapterheadstartvskip {%
%	\vspace*{1\baselineskip plus .1\baselineskip minus .067\baselineskip}}

% hint boxes
\newmdenv[%
roundcorner=5pt,
backgroundcolor=gray!10!white,
linewidth=1pt,
%innertopmargin=0pt,
%skipabove=0pt,
linecolor=gray!75,
font={\normalfont\sffamily},
%frametitle={Hint},
%frametitlerule=true,
%frametitlebackgroundcolor=yellow,
%frametitlefont={\normalfont\sffamily},
%frametitleaboveskip=8pt,
%frametitlebelowskip=8pt,
subtitlebackgroundcolor=gray!40,
subtitlefont={\normalfont\sffamily},
]{hint}

\newmdenv[%
roundcorner=5pt,
subtitlebelowline=true,subtitleaboveline=true,
subtitlebackgroundcolor=yellow!70!white,
backgroundcolor=blue!20!white,
frametitle={Theorem},
frametitlerule=true,
frametitlebackgroundcolor=yellow!70!white,
]{subtitleenv}

\newcommand{\rcg}{\rowcolor{gray!40}}
\newcommand{\rcgg}{\rowcolor{gray!10}}

 %\overfullrule=5pt
 %\underfullrule=5pt
 
 
\RedeclareSectionCommand[
	beforeskip=-.75sp,
	afterskip=1.5\baselineskip]{chapter}
\RedeclareSectionCommand[
	beforeskip=1.0\baselineskip,
	afterskip=.75\baselineskip]{section}
\RedeclareSectionCommand[
	beforeskip=-.75\baselineskip,
	afterskip=.5\baselineskip]{subsection}
\RedeclareSectionCommand[
	beforeskip=-.75\baselineskip,
	afterskip=.25\baselineskip]{subsubsection}
\RedeclareSectionCommand[
	beforeskip=.75\baselineskip,
	afterskip=.25\baselineskip]{paragraph}

%\setlength{\footheight}{56pt}
%\setlength{\headheight}{44.83598pt}
